% Generated from raw_evidence in inputs.json.
\begingroup\footnotesize
\begin{longtable}{@{}>{\raggedright\arraybackslash}p{25mm}>{\raggedright\arraybackslash}p{61mm}>{\raggedright\arraybackslash}p{70mm}@{}}
\caption{紧凑证据表：原值与本分析采用方式。PDF定位为本地页序。}\\
\toprule 来源/定位 & 原值与条件 & 采用、换算及限制\\\midrule\endfirsthead
\toprule 来源/定位 & 原值与条件 & 采用、换算及限制\\\midrule\endhead
NAND-04 PDF p.2 II-A, Fig.1; p.3 II-C, Fig.3 & 13824 BL $\times$ 32 WL $\times$ 3 SSL/block; 64 blocks/subarray; BL length 144 $\mu$m。作者CIM结构/RC模型，32 nm逻辑外围；不是量产页手册 & 原样保留物理几何。两subarray、128 blocks构成一份矩阵；一WL/一SSL的一页为13824 bit=1728 Byte。\\
NAND-05 PDF p.1 II-(b); p.3 Figs.5--6 & 16-layer, 64 Gb SLC; ION mean 2 nA, sigma 0.3 nA; IOFF <0.5 pA; Vg=1 V, VBL=0.2 V, Vpass=4.5 V。SGVC实际器件电流；full-block random code；非完整CIM芯片速度 & 二阈值SLC依据。NAND-04的32 WL结构另属模型，不把16层和32层说成同一实测芯片。\\
NAND-04 PDF p.3 Table 1及II-C & WL setup 303 ns; BL setup 12 ns @50\% sparsity; SL setup 530--750 ns @max-bit input; SL capacitance $\sim$16 pF。RC/HSPICE估计，64-block subarray；SL不是完整SLC read cycle & WL每输出组一次(8次/向量)；BL和SL每求值一次(64次)。SL中点640 ns仅插值情景；复制108份使电流接近原负载，范围仍以适配成功为条件。\\
NAND-05 PDF p.1 II-(a); p.4 Fig.15 & Tread $\sim$1 $\mu$s with multibit output; $\sim$10 $\mu$A MAC current; 7--8-bit SA resolution。加速器设计预估，非单cell测量；包括多位SA输出 & 仅对530--750 ns阵列建立加共同ADC的微秒量级交叉检查，不把1 $\mu$s再叠加或当独立硅证据。\\
NAND-04 PDF p.2 II-A; p.4 II-C & Both states programmed with write-verify; programming voltage $\sim$20 V。两阈值都细调以控制导通电流 & 完整P必须含charge-pump建立、所有program/verify循环、恢复及可计算终点；完整P未给，保持未知。\\
NAND-05 PDF p.2 II-(e); p.4 Fig.12 & Sacrifice one or two WLs per block for known-code SA calibration; on-the-fly calibration。读取已知码后校准SA；重写/漂移可能需要重新校准 & 主情景1个校准WL，另核对2个。擦除后每块增加1或2个完整校准页program；C\_\allowbreak{}cal为整矩阵事务的总校准时间，未报告，参数化。\\
NAND-01 PDF p.2 Key Specifications & MLC: tPROG 1300/2500 $\mu$s, tBERS 15/45 ms, tR 77 $\mu$s; TLC: 1630/5000 $\mu$s, 15/45 ms, 100 $\mu$s (typ/max)。MLC 2b/c与TLC 3b/c分列；16,384+2,208 Byte/page; 1024/1536 pages/block & 仅说明商品存储模式完整周期及粒度不同；全部不代入SLC主情景，也不充当SLC上/下界。\\
NAND-02 PDF p.1及p.3 Fig.15.1.5 & Fast SLC burst $\sim$11\% interior tPROG reduction; first/last $\sim$5.5\%; pump held between pages。332-WL QLC产品的SLC-cache burst；未给SLC绝对P或E & 证明泵建立/最终恢复必须计入有限窗口；不把11\%或QLC >85 MB/s作为SLC绝对时间，不采用burst收益。\\
NAND-03 PDF p.1及p.2 Fig.30.5.6 & QLC 4b/c; 6 planes; 16 kB page; tPROG 1.3 ms; die throughput 75 MB/s。321-layer 2 Tb芯片；6-plane并行 & 仅核对die/plane吞吐与页cycle不同；不代入SLC单更新域。\\
公共基线：参数JSON与R0 & 128$\times$128 INT8; BS=128 Byte; BR(full)=16384 Byte; R0: 8 weight planes, 16 ADC/plane; (TI,TA,TD)=(2,10,2)/(5,20,5)/(10,50,10) ns; 128-bit write interface。已接受公共工程选择；ACIM校准近似部分和、128$\times$24-bit输出 & 逻辑合同和时隙不变；所有次数由导入的acim\_\allowbreak{}service/front\_\allowbreak{}ns/program\_\allowbreak{}sequence\_\allowbreak{}ns接口计算。\\
\bottomrule\end{longtable}\endgroup
